% Peekaboo-VL paper: 09_conclusion. Included from ../peekaboo_report.tex.
\section{Conclusion}
We set out to give a video language model the camera motion it cannot see
well, by measuring it and writing it down instead of learning it. A
few-millisecond CPU probe that describes the motion in the answers' own
words, plus LoRA on an off-the-shelf video language model, scores 3.5
points above a purpose-built motion architecture on ReVA, with no new
network parts and one GPU. Most of that margin comes from the newer
backbone and from fine-tuning on the seen videos. The fine-tuned model
does read the hint, but the hint itself adds little.

Three things we learned should hold beyond this benchmark. The test regime
is seen-video, so a development split without training-answer leaks is
essential. Text alone solves 56\%, so the video's real contribution has to
be measured against a blind baseline. And a hint's wording matters more
than the measurement behind it. Temporal grounding remains the weakest
task: training with 32 shifted frames and a drawn clock raised it by 2.7
points, but its errors are still near-misses of about one timestamp. We
hope this encourages more video QA benchmarks to report text-only
baselines and per-task numbers.

The code, the training recipe and the released LoRA adapter are available
at \url{https://github.com/LakoreAI/peekaboo-vl} and
\url{https://huggingface.co/LakoreAI/peekaboo-vl}.
